% =============================================================================
% ISI Paper 3 -- Analytical Brief (English)
% External Supplier Concentration in the EU-27
% Sections 1--7
% =============================================================================

% ─────────────────────────────────────────────────────────────────────────────
\section{Executive Summary}
\label{sec:summary}
% ─────────────────────────────────────────────────────────────────────────────

\begin{sloppypar}
\textbf{Core finding.}
The International Sovereignty Index (\ISI) is a six-dimensional
composite indicator of external supplier concentration.  Its
empirical application to the EU-27 (vintage~2024) shows that the
cross-sectional variance of composite scores is predominantly
determined by two of the six axes: defence and logistics.  This
result is a descriptive variance finding --- not a conceptual
reduction of the index to two dimensions.
\end{sloppypar}

\medskip
\textbf{Principal empirical results:}
\begin{itemize}[leftmargin=2em]
  \item The \ISI composite score ranges from France~(0.236;
    rank~27) to Malta~(0.517; rank~1).  The EU-27 mean
    is~0.344 with a standard deviation of~0.070.

  \item \begin{sloppypar}The covariance-based variance
    decomposition yields the following: the defence axis
    contributes~35.1\,\% to composite variance; the logistics
    axis~34.3\,\%.  Together, these two axes account
    for~69.4\,\% of total cross-sectional
    variance.\end{sloppypar}

  \item The remaining four axes --- finance, energy, technology,
    critical inputs --- jointly contribute~30.6\,\%.  Their
    values are distributed comparatively homogeneously across
    the EU-27 cross-section.

  \item In 16 of 27~member states, the defence axis records the
    highest individual axis score; in the remaining~11, the
    logistics axis does.  None of the other four axes records
    the highest score in any member state.

  \item Removing the defence axis from the composite produces
    the largest rank displacement
    (Spearman~$\rho = 0.833$ relative to the baseline ranking).
    All other single-axis removals yield correlations
    above~0.94.

  \item Under Dirichlet weight perturbation (10\,000~draws),
    the mean rank-order correlation is
    $\rho = 0.979$.  The rank ordering is stable under
    moderate weight variations.

  \item \textbf{Austria} ranks~8th (composite score~0.383;
    moderately concentrated).  Highest axis score:
    defence~(0.585), followed by logistics~(0.526).  The
    rank position is structurally stable --- under no
    single-axis perturbation does Austria shift by more
    than two ranks.
\end{itemize}


\newpage
% ─────────────────────────────────────────────────────────────────────────────
\section{What the ISI Measures}
\label{sec:what-the-isi-measures}
% ─────────────────────────────────────────────────────────────────────────────

The \ISI is a composite indicator that measures the
\emph{concentration of external suppliers}.  For each EU member
state, it captures the extent to which imports in six strategically
relevant domains depend on a small number of external suppliers.
The index comprises six equally weighted dimensions:

\begin{table}[htbp]
\centering\small
\caption{The six \ISI axes and their coverage domains.}
\label{tab:axes}
\begin{tabular}{@{}lp{0.64\linewidth}@{}}
\toprule
\textbf{Axis} & \textbf{Coverage domain} \\
\midrule
Finance &
  Cross-border banking and portfolio
  concentration \\
Energy &
  Dependence on a small number of primary
  energy suppliers \\
Technology &
  Import concentration in
  high-technology goods \\
Defence &
  Supplier concentration in arms transfers and
  defence-relevant imports \\
Critical Inputs &
  Concentration in critical raw materials and
  intermediate goods \\
Logistics &
  Dependence on a small number of transport
  corridors and shipping routes \\
\bottomrule
\end{tabular}
\end{table}

Each axis score is based on a Herfindahl-Hirschman-type
concentration measure, normalised to the unit interval
$[0,\,1]$.  A value of zero denotes complete
diversification; a value near one denotes near-complete
dependence on a single external supplier.  All six
axes are identically scaled.

The composite score is the \emph{equal-weighted arithmetic
mean} of the six axis scores.  The aggregation is
deterministic --- no estimation, no latent factors,
no model selection.  The methodology is fully documented:
Paper~1~\cite{isi_paper1} presents the methodological
framework; Paper~2~\cite{isi_paper2} documents the empirical
application to the EU-27.

\textbf{Scope and interpretation.}
The \ISI does not measure ``sovereignty'' in the normative or
constitutional sense.  It captures a specific, quantifiable
structural property: the degree of concentration of external
supply relationships.  High concentration is not inherently
negative --- it may reflect geographical proximity, trade
efficiency, or alliance structures.  It does, however, imply
that the country's position is exposed to disruptions in a
small number of supply relationships.


\newpage
% ─────────────────────────────────────────────────────────────────────────────
\section{Structural Differentiation within the EU-27}
\label{sec:differentiation}
% ─────────────────────────────────────────────────────────────────────────────

\subsection{Variance Contributions of the Six Axes}

The \ISI is conceptually a six-dimensional index.  All six axes
enter the composite score with equal weight and contribute to the
absolute concentration level of each country.  The empirical
question of this section is different: which axes determine the
\emph{differences} among member states?

Since all six axes are identically normalised to the
interval~$[0,\,1]$ and enter the composite with equal weight, the
covariance-based variance decomposition is directly interpretable.
The cross-sectional variance of composite scores is predominantly
determined by two axes:

\begin{table}[htbp]
\centering\small
\caption{Variance decomposition of the \ISI composite score, EU-27.}
\label{tab:variance}
\begin{tabular}{@{}lr@{}}
\toprule
\textbf{Axis} & \textbf{Share of composite variance} \\
\midrule
Defence            & 35.1\,\% \\
Logistics          & 34.3\,\% \\
Critical Inputs    & 20.9\,\% \\
Technology         &  5.2\,\% \\
Energy             &  3.5\,\% \\
Finance            &  1.0\,\% \\
\midrule
Defence + Logistics      & 69.4\,\% \\
Remaining four axes      & 30.6\,\% \\
\bottomrule
\end{tabular}
\end{table}

Approximately 64\,\% of total variance is attributable to
own-variance (dispersion of each axis across countries); 36\,\%
is attributable to cross-covariances (joint variation across
axes).  Defence and logistics dominate both components.

\textbf{Methodological note.}
The variance decomposition is purely descriptive.  It identifies
which axes statistically determine cross-sectional variance --- not
which axes are politically paramount.  An axis may be of the
highest political relevance yet contribute little to cross-sectional
differentiation if all member states exhibit similar values on that
axis.  Political priority and statistical differentiation are
analytically distinct categories.

The finding implies that structural heterogeneity within the EU is
not uniformly distributed across all strategic domains but
concentrated in a small number of dimensions.  Integration-policy
tensions are more likely where national profiles diverge
substantially.

\subsection{Why Energy and Technology Differentiate Less}

Energy dependence is of high relevance for many member states.
Across the EU-27 cross-section, however, it differentiates
countries less strongly than defence and logistics.  The reason is
empirical: energy concentration scores exhibit comparatively low
dispersion.  Most member states display similar degrees of
energy supplier concentration.  The same applies to the technology
axis.  Both axes contribute to the absolute concentration level of
each country but generate smaller differences among countries than
defence and logistics.

\subsection{Axis Dominance at the Country Level}

In 16~member states, the defence axis records the highest
individual axis score; in~11, the logistics axis does.  None of the
remaining four axes records the highest score in any member state.
This pattern describes the empirical distribution of axis
maxima --- it does not imply that the \ISI can be conceptually
reduced to two axes.  The index remains a six-dimensional
instrument; the observed two-axis dominance is a finding about the
data, not about the construction.


% ─────────────────────────────────────────────────────────────────────────────
\section{Austria in the EU-27 Context}
\label{sec:austria}
% ─────────────────────────────────────────────────────────────────────────────

Austria ranks~8th of~27 with a composite score of~0.383.  It is
classified as ``moderately concentrated'' --- the same category as
23 of the 27~member states.

\begin{table}[htbp]
\centering\small
\caption{Austria's axis profile.}
\label{tab:austria}
\begin{tabular}{@{}lrr@{}}
\toprule
\textbf{Axis} & \textbf{Score} & \textbf{Axis rank (of~27)} \\
\midrule
Defence            & 0.585 & 11th \\
Logistics          & 0.526 & 10th \\
Energy             & 0.456 &  3rd \\
Critical Inputs    & 0.319 &  7th \\
Technology         & 0.261 &  5th \\
Finance            & 0.149 & 11th \\
\bottomrule
\end{tabular}
\end{table}

Austria's highest axis score is on defence~(0.585), followed by
logistics~(0.526).  The gap between the two is~0.059 --- Austria
does not exhibit a strongly asymmetric axis profile.

\textbf{Energy concentration.}
Austria's third-highest axis score is energy~(0.456), rank~3 in
the EU.  This reflects the documented dependence on a small number
of gas suppliers.  Because energy concentration is distributed
comparatively homogeneously across the EU cross-section, it
affects Austria's composite ranking only marginally.

\textbf{Profile balance.}
Austria's coefficient of variation across the six axes is~0.40 ---
a comparatively balanced axis profile.  For comparison:
Sweden (rank~7; composite score~0.399) records a defence score
of~0.881 but a finance score of only~0.116 --- a substantially
more asymmetric profile.

\textbf{Rank stability.}
Under leave-one-out analysis --- the systematic removal of each
individual axis --- Austria's rank shifts by at most two positions.
Under Dirichlet weight perturbation (10\,000~Monte Carlo draws),
Austria's rank standard deviation is approximately~1.1.  The
95\,\% confidence interval spans ranks~6 to~10.  Austria's
position in the upper middle range is structurally robust.

\textbf{Neighbourhood.}
Austria is situated between Sweden (rank~7; 0.399) and
Italy (rank~9; 0.374).  The distances are~0.016 upward
and~0.009 downward.  Moderate changes to a single axis could
shift Austria by one to two rank positions; a move into the top~5
or below rank~12 would require substantial structural changes.


% ─────────────────────────────────────────────────────────────────────────────
\section{Robustness of Results}
\label{sec:robustness}
% ─────────────────────────────────────────────────────────────────────────────

A legitimate concern regarding any composite indicator is the
dependence of results on specific methodological choices.  Three
tests address this concern.

\subsection{Leave-one-out Analysis}

When each axis is individually removed from the composite, six
alternative rankings result.  Removing the defence axis produces
the lowest correlation with the baseline ranking
(Spearman~$\rho = 0.833$) --- confirming that defence is the
single most influential axis for rank differentiation.
Removing logistics yields~$\rho = 0.926$.  All remaining
removals produce correlations above~0.94.

The ordering is monotonic: axes with higher variance contributions
produce larger rank displacements upon removal.

\subsection{Weight Stability}

Equal weighting is a design choice.  To test its consequences,
10\,000~alternative weight vectors are drawn from a symmetric
Dirichlet distribution and the composite score is recomputed for
each draw.  The mean rank-order correlation across all draws
is~0.979 --- the rank ordering is stable under moderate weight
variations.

The highest rank volatility is exhibited by the Netherlands
(standard deviation $\approx$~3.07~rank positions).  This
volatility reflects the asymmetric axis profile of the
Netherlands: the defence score~(0.960) lies far above the level
of the remaining five axes.  Stronger weighting of the defence
axis raises the Dutch rank position; weaker weighting lowers it.
For the large majority of member states, the ranking is stable to
within one to two positions.

\subsection{Correlation Structure}

The six axes are largely independent of one another.  Of
15~pairwise correlations, only two are statistically significant
after Bonferroni correction: energy--technology and critical
inputs--logistics.  The finance axis is nearly orthogonal to all
other axes (highest $|t|$-value: 1.11).

This independence is analytically advantageous: the variance
decomposition reflects genuine structural content rather than
statistical redundancy.  The low multicollinearity prevents
artificial variance inflation of individual axes.


% ─────────────────────────────────────────────────────────────────────────────
\section{Interpretation and Analytical Implications}
\label{sec:implications}
% ─────────────────────────────────────────────────────────────────────────────

The \ISI does not prescribe policy.  It provides an empirical
measurement basis.

\textbf{Cross-sectional variance and absolute concentration.}
Defence and logistics determine cross-sectional variance.  It does
\emph{not} follow that the remaining axes are irrelevant for
individual member states.  A country may exhibit a high absolute
concentration score on the energy axis without that score
differentiating countries in the cross-section.  Political priority
and statistical differentiation are analytically distinct
categories.

\textbf{Empirical grounding of strategic debates.}
The term ``strategic autonomy'' is ubiquitous in EU policy
discourse but rarely linked to comparable, quantitative data.  The
\ISI provides such an anchor point.  The variance decomposition
shows that defence procurement and logistics corridors generate the
largest differences among member states --- domains that receive
comparatively little attention in public debate relative to energy
and technology.

\begin{sloppypar}
\textbf{Energy and technology.}
Energy dependence is of high importance for numerous member states.
Across the EU-27 cross-section, however, it differentiates
countries less strongly than defence and logistics because energy
concentration scores are distributed comparatively uniformly.  The
same applies to technology.  Both axes require independent
analysis; their low contribution to cross-sectional dispersion does
not diminish their substantive significance.
\end{sloppypar}

\textbf{Austria.}
Rank~8 places Austria in the upper middle range of external
concentration.  The comparatively balanced axis profile exhibits
neither a singular concentration peak nor an exceptional
diversification strength.

\textbf{Measurement, not evaluation.}
A high composite score does not mean that a country ``performs
poorly.''  Malta's rank~1 reflects size and geographical location,
not policy failure.  France's rank~27 reflects the structural
diversification of a large, well-connected economy.  The index
measures concentration --- what policymakers derive from it is a
separate question.


% ─────────────────────────────────────────────────────────────────────────────
\section{Limitations}
\label{sec:limitations}
% ─────────────────────────────────────────────────────────────────────────────

\textbf{Re-export bias.}
Trade statistics capture bilateral flows but do not always identify
the actual origin of goods.  Countries that function as trade
hubs (such as the Netherlands or Belgium) may appear more
diversified than they actually are.

\textbf{Amplification effect for small economies.}
Small, open economies with a limited number of trading partners
mechanically generate higher concentration scores.  Malta and
Cyprus rank~1st and~2nd in part because geography and size
constrain their sourcing options.  This is a real feature of their
situation, not a measurement error --- it should, however, be
interpreted accordingly.

\textbf{Cross-sectional limitation.}
The dataset captures the 2024 vintage.  Supplier concentration may
change as a result of political decisions, market developments, or
geopolitical events.  A single cross-section cannot capture
dynamics.  Future \ISI vintages will enable trend analysis.

\textbf{Equal weighting.}
The aggregation treats all six axes as equally important.  This is
a normative assumption.  A decision-maker who regards energy
dependence as paramount would apply different weights and obtain a
different ranking.  The Dirichlet perturbation analysis shows that
moderate deviations from equal weighting do not fundamentally alter
the results.

\vfill
\begin{center}
\rule{0.4\linewidth}{0.4pt}
\end{center}
\smallskip
\noindent\textit{This brief is based on the ISI Paper Series
(Papers~1--2), prepared by the Research Unit of the International
Sovereignty Institute.  The underlying data and the complete
methodology are publicly documented.}
